\section{Hilbert comparisons with specified domains}
\label{v4-arith:mc:comparisons}

The continuous Mellin unitary and the global arithmetic character
act on different spaces. A comparison begins with the inner products
and densely defined maps between these spaces. All Hilbert spaces
below are complex, with inner products linear in the first variable.

\begin{proposition}[extension of a core isometry]
\label{v4-arith:mc:core-unitary}
Let $T_j$ be densely defined operators on Hilbert spaces
$\mathcal H_j$, with domains $\mathcal C_j$, for $j=1,2$.
Suppose a complex-linear bijection $U_0:\mathcal C_1\to\mathcal C_2$
satisfies
\[
 \langle U_0v,U_0w\rangle_2=\langle v,w\rangle_1,
 \qquad U_0T_1v=T_2U_0v,
\]
where $T_1\mathcal C_1\subseteq\mathcal C_1$ and
$T_2\mathcal C_2\subseteq\mathcal C_2$.
Then $U_0$ extends uniquely to a unitary $U:\mathcal H_1\to\mathcal H_2$.
It identifies the closures of the graphs and the adjoint domains:
\[
 U\operatorname{Dom}T_1^*=\operatorname{Dom}T_2^*,\qquad
 UT_1^*=T_2^*U.
\]
If the operators are closable, then
$U\overline T_1=\overline T_2U$, with equality of the transported
domains. In particular, essential self-adjointness and deficiency
spaces are preserved.
\end{proposition}
\begin{proof}
For $v\in\mathcal H_1$, choose $v_n\in\mathcal C_1$ with $v_n\to v$
and define $Uv=\lim U_0v_n$. The isometry identity proves existence
and independence of the sequence. The range is closed and contains
the dense space $\mathcal C_2$, so $U$ is onto.
The unitary $U\oplus U$ maps the graph of $T_1$ onto that of $T_2$.
It therefore maps their graph closures onto one another.
For $y\in\mathcal H_1$, the functional
$v\mapsto\langle T_1v,y\rangle_1$ is bounded in the Hilbert norm
exactly when $U_0v\mapsto\langle T_2U_0v,Uy\rangle_2$ is bounded.
This gives the adjoint-domain identity and the adjoint action.
Hence $U\ker(T_1^*\mp i)=\ker(T_2^*\mp i)$.
\end{proof}

The same proof applies without invariance of the cores when a
unitary $U$ is already given and $UT_1=T_2U$ on the exact domains.
In particular, if $H_j$ are self-adjoint, $\mathcal C_j$ are their
graph cores, and the preceding hypotheses hold for their restrictions,
then $U\operatorname{Dom}H_1=\operatorname{Dom}H_2$ and $UH_1=H_2U$.

\begin{proposition}[transport of a complete spectral realization]
\label{v4-arith:mc:transport}
Let $H_j$ be self-adjoint on $\mathcal H_j$ and let
$U:\mathcal H_1\to\mathcal H_2$ be unitary, with
$U\operatorname{Dom}H_1=\operatorname{Dom}H_2$ and $UH_1=H_2U$.
Let $A_1:\operatorname{Dom}A_1\subset\mathcal H_{\mathrm{BC}}
\to\mathcal H_1$ be a closed densely defined map. Define
$A_2=UA_1$ on the same domain.
Then $A_2$ is closed. Injectivity and dense range are equivalent
for $A_1$ and $A_2$. On any specified $N_{\mathrm{occ}}$-invariant
subspace $\mathcal C_N\subseteq\operatorname{Dom}A_1\cap
\operatorname{Dom}N_{\mathrm{occ}}$ with
$A_1\mathcal C_N\subseteq\operatorname{Dom}H_1$, one has
\[
 A_1N_{\mathrm{occ}}=H_1A_1
 \quad\Longleftrightarrow\quad
 A_2N_{\mathrm{occ}}=H_2A_2.
\]
The shifted operators $\Theta_j=1/2+iH_j$ have the same eigenvalues,
eigenspace dimensions, and spectral type.
For every bounded Borel $\varphi$ with $\varphi(H_1)$ trace-class,
\[
 \operatorname{Tr}\varphi(H_2)=\operatorname{Tr}\varphi(H_1).
\]
\end{proposition}
\begin{proof}
The unitary $1\oplus U$ maps the graph of $A_1$ onto that of $A_2$,
so closedness, kernel, and closure of the range are preserved.
Domain compatibility and $UH_1=H_2U$ give the displayed equivalence
by multiplication with $U$ or $U^{-1}$.
For nonreal $z$, the resolvents satisfy
$(H_2-z)^{-1}=U(H_1-z)^{-1}U^{-1}$.
Uniqueness of the spectral resolution then gives
$\varphi(H_2)=U\varphi(H_1)U^{-1}$ and transports spectral subspaces.
Taking an orthonormal basis and its image under $U$ proves trace
invariance whenever the trace is absolutely defined.
\end{proof}

\begin{corollary}[a common determinant normalization]
\label{v4-arith:mc:determinant}
In Proposition~\ref{v4-arith:mc:transport}, suppose that on a fixed
domain of $(s,u)$ the powers $(s-\Theta_j)^{-u}$ use the same
holomorphic branch and are trace-class. Suppose their traces admit
meromorphic continuation to a neighborhood of $u=0$, regular there.
Then the determinants
\[
 d_j(s)=\exp\left(-\left.\partial_u
          \operatorname{Tr}(s-\Theta_j)^{-u}\right|_{u=0}\right)
\]
are equal. Any identical polar-factor removal on the two sides
preserves this equality. Thus $d_1=\xi_{\mathbb R}$ transfers to
$d_2=\xi_{\mathbb R}$ with the same normalization.
\end{corollary}
\begin{proof}
Functional calculus and trace invariance identify the traces on
their initial domain. Uniqueness of meromorphic continuation
identifies their derivatives at zero. Exponentiation proves the claim.
Neither existence of the initial trace-class powers nor continuation
follows from an identity of formal eigenvalue lists.
\end{proof}

\begin{definition}[the Fourier test algebra]
\label{v4-arith:mc:test-algebra}
For $f\in\mathscr A=C_c^\infty(\mathbb R)$ define the entire function
$\widehat f(z)=\int_{\mathbb R}f(t)e^{izt}dt$.
Let $\mathscr P=\{\widehat f:f\in\mathscr A\}$, with the topology
transported from compact smooth tests by this injective Fourier map.
Its involution is $h^\star(z)=\overline{h(\overline z)}$.
For a specified continuous arithmetic functional
$\mathfrak W:\mathscr P\to\mathbb C$, put
$W(f)=\mathfrak W(\widehat f)$.
\end{definition}

Fubini on compact supports gives
$\widehat{f*g}=\widehat f\,\widehat g$ and
$\widehat{g^*}=(\widehat g)^\star$, where
$g^*(t)=\overline{g(-t)}$.
Injectivity follows from the Fourier unitary of
Lemma~\ref{mo:mellin-unitary} at $c=0$.
Thus $W(f*g^*)=\mathfrak W(\widehat f(\widehat g)^\star)$
is an exact comparison of the two test carriers.
If $\mathfrak W$ is specified on a larger Paley--Wiener class,
extending an identity from $\mathscr P$ to that class requires
the stated density and continuity hypotheses.

\begin{proposition}[a Hilbert space from a positive convolution form]
\label{v4-arith:mc:positive-form}
Let $\mathscr A=C_c^\infty(\mathbb R)$, with additive convolution
and $f^*(t)=\overline{f(-t)}$. Let $W:\mathscr A\to\mathbb C$
be a continuous linear functional such that
$q_W(f,g)=W(f*g^*)$ is Hermitian and $q_W(f,f)\geq0$.
Its nullspace $\mathscr N$ is a linear subspace, and
$\mathcal H_W=\overline{\mathscr A/\mathscr N}^{q_W}$ is a Hilbert space.
The maps $V_u[f]=[f(\,\cdot+u)]$ extend to a strongly continuous
unitary group on $\mathcal H_W$.
For $h\in\mathscr A$, the convolution action
$\pi_W(h)[f]=[h*f]$ extends to a bounded operator with
$\|\pi_W(h)\|\leq\|h\|_{L^1}$ and
$\pi_W(h^*)=\pi_W(h)^*$.
If a linear map $C:\mathscr A\to\mathcal K$ has dense range and
$\langle Cf,Cg\rangle_{\mathcal K}=q_W(f,g)$, it induces a unique
unitary $\mathcal H_W\to\mathcal K$.
\end{proposition}
\begin{proof}
Positivity gives Cauchy--Schwarz for $q_W$ by applying it to
$q_W(f+zg,f+zg)$ for all $z\in\mathbb C$.
Thus null vectors are orthogonal to every test, proving the quotient
assertion. Translation by $u$ in both arguments leaves $f*g^*$
unchanged. Hence $V_u$ is an isometry with inverse $V_{-u}$.
As $u\to0$, translated tests converge in the compact-smooth topology
on a common compact support. Continuity of convolution and $W$
gives $\|V_u[f]-[f]\|\to0$. Density and the isometry extend strong
continuity to the completion. The last claim follows by completing
the isometry $[f]\mapsto Cf$.
For convolution, integrate the strongly continuous unitaries:
\[
 \pi_W(h)v=\int h(u)V_{-u}v\,du.
\]
Riemann sums and the triangle inequality give the displayed norm
bound. On test vectors, the same sums converge in the compact-smooth
topology to $h*f$, so continuity of the quotient map identifies
the integral with the asserted action. Taking adjoints of the
integral and substituting $u\mapsto-u$ proves the involution identity.
Fubini proves $\pi_W(h_1*h_2)=\pi_W(h_1)\pi_W(h_2)$.
\end{proof}

This construction supplies a positive Hilbert space only when the
specified arithmetic functional satisfies the positivity hypothesis.
It supplies neither the occupation map nor an ordinary character trace.
For example, if $W(h)=r\int h(t)e^{i\gamma t}dt$ with integer $r>1$,
the quotient has dimension one: its norm is
$r|\int f(t)e^{i\gamma t}dt|^2$.
An ordinary representation with trace multiplicity $r$ at $\gamma$
has dimension $r$. A positive cyclic form therefore does not recover
ordinary trace multiplicity by itself.

\begin{proposition}[the ranges of scattering isometries]
\label{v4-arith:mc:scattering}
Let $\Omega_\pm:\mathcal K_0\to\mathcal K$ be isometries of
complex Hilbert spaces and put $S=\Omega_+^*\Omega_-$.
Then $S$ is unitary exactly when
$\operatorname{Ran}\Omega_+=\operatorname{Ran}\Omega_-$.
\end{proposition}
\begin{proof}
Write $P_+=\Omega_+\Omega_+^*$ for the orthogonal projection onto
$\operatorname{Ran}\Omega_+$. If $S$ is an isometry, then
\[
 \|(1-P_+)\Omega_-v\|^2
 =\|v\|^2-\|Sv\|^2=0.
\]
Thus the minus range is contained in the plus range.
If $S$ is unitary, the same argument for $S^*$ gives the reverse
inclusion. Conversely, equal ranges give
$S^*S=1=SS^*$ by inserting their common range projection.
\end{proof}

Equal scattering ranges do not determine an arithmetic determinant.
For the constant dynamics on the one-dimensional space $\mathbb C$,
take $\Omega_+=\Omega_-=1$. Then $S=1$, while the shifted operator
has the eigenvalue $1/2$, which is excluded from the actual xi divisor
by Proposition~\ref{mo:vacuum}.
